\section{Certificate schemas and executable entry points}

\subsection{Canonical scalar and polynomial encodings}

A rational is encoded as the integer pair \code{[numerator, denominator]}, with
positive denominator and greatest common divisor one. Python booleans are
rejected even though they are subclasses of integers. Floats are rejected.
A sparse polynomial in $D$ variables is a strictly lexicographically ordered
list
\begin{lstlisting}
[[[e1, ..., eD], [numerator, denominator]], ...]
\end{lstlisting}
with nonnegative exponents, no repeated exponent vector, and no explicit zero
coefficient. The empty list denotes the zero polynomial. Constant polynomials
have the all-zero exponent vector. The source problem establishes the variable
order; the certificate cannot reorder it.

The checker uses exponent-vector addition for monomial multiplication and
collects equal exponents by exact addition. Substitution is evaluated by repeated
squaring and cached powers. This implementation is independent of SymPy's
internal polynomial classes. Both search and checker still share the published
mathematical schema, so common specification mistakes remain possible; the
paper's derivations and the concrete differential controls help expose them but
do not formally exclude them.

\subsection{Linear and ideal payloads}

The linear problem supplies \code{dimension}, \code{actions}, \code{initial},
and \code{target}. Its certificate supplies \code{basis},
\code{target\_weights}, and \code{action\_weights}. Matrices are row-major.
Every action must have one correctly dimensioned coefficient matrix.

The ideal problem additionally supplies \code{parameters}, with its initial
coordinate polynomials in those parameter variables. Its certificate supplies
\code{generators}, \code{target\_weights}, and nested
\code{action\_weights}. The latter is indexed by action, destination generator,
and source generator. Each entry is a polynomial in the state variables. There
is no optional ``trust this Gr\"obner basis'' field.

A concrete counterexample payload contains a forward input \code{word} and its
nonzero rational \code{value}; polynomial cases also contain concrete integer
\code{parameters}. Replaying the original problem must reproduce that value.
A nonzero value alone, without a valid word and initialization, is not accepted.

\subsection{Summation payloads}

The summation problem supplies \code{power} and a polynomial \code{weight} in
$k$. The certificate supplies \code{p0}, \code{p1}, \code{r}, and
\code{relation\_multiplier}. Their variable orders are respectively $(n)$,
$(n)$, $(n,k)$, and $(A,B,n,k)$. The local checker validates
\eqref{eq:sum-cert}. Endpoints are justified by the fixed shape theorem, not by
certificate-supplied numerical endpoint claims. An extension that changes this
shape must provide a new boundary theorem.

\Needspace{23\baselineskip}
\subsection{Actual use from Python}

The following is executable against the delivered package:
\begin{lstlisting}[language=Python]
import sympy as s
from prototype.search import ideal_problem, ideal_search
from prototype.checkers import verify

x, y, z, w = s.symbols("x y z w")
a, b = s.symbols("a b")
problem = ideal_problem(
    (x, y, z, w),
    [(y*y, x, w*w, z)],
    (a, b),
    [a, b, a, -b],
    x - z,
)
result = ideal_search(problem, max_extensions=12)
assert result["status"] == "certified"
assert verify(problem, result["certificate"])
\end{lstlisting}

A user may replace the target or updates, but the certificate must then be
rechecked against that changed problem. Reusing a previous result object or
matching a case name does not establish validity.

\section{Reproduction and archive map}

From the archive root:
\begin{lstlisting}
python -m pip install -r requirements.txt
python -m prototype.run_experiments --output reproduced-results
python -S prototype/checkers.py reproduced-results/certificates \
    reproduced-results/counterexamples
python -m pytest -q
\end{lstlisting}
The regression tests replay the retained \code{results/} corpus; the generation
command writes a distinct directory by default. Timings will vary across
machines. The deterministic outcome categories and exact identities, not the
millisecond values, are the reproducibility targets.

To build the article, run \code{python build.py}. It invokes \code{pdflatex}
three times in a temporary output directory and copies the final PDF into
\code{article/}. The complete source uses ordinary LaTeX packages; no font files
are included. The generated \code{metrics.tex} and \code{recurrences.tex} reflect
the retained run, not an automatically substituted local timing run.

The \code{lean/} directory is a separate, uncompiled specimen project. On a
machine with the pinned Lean toolchain and fetched Mathlib dependencies,
\code{lake build} is the intended first validation command. Its successful
execution is \emph{not} part of the recorded evidence supplied with this article.

\begin{longtable}{@{}p{0.34\linewidth}p{0.58\linewidth}@{}}
\toprule
Path & Contents\\
\midrule\endhead
\code{article/} & Main TeX, section sources, generated tables, references, and rendered PDF.\\
\code{prototype/search.py} & The three untrusted exact search algorithms.\\
\code{prototype/checkers.py} & Standard-library proof and counterexample replay.\\
\code{prototype/run\_experiments.py} & Deterministic corpus, timing protocol, mutation generation, and exports.\\
\code{tests/} & Regression, schema, boundary, and differential tests.\\
\code{results/certificates/} & 38 compact problem-and-certificate bundles.\\
\code{results/counterexamples/} & 17 original-problem concrete counterexamples.\\
\code{results/summary.json} & Per-case outcomes, settings, timing statistics, and environment.\\
\code{results/mutations.json} & All 818 single-entry perturbation outcomes.\\
\code{results/*.log} & Accepted runs and explicitly separated development diagnostics.\\
\code{docs/} & Repository audit, validation scope, and source manifest.\\
\code{lean/} & Uncompiled semantic induction and worked replay specimens.\\
\bottomrule
\end{longtable}

\section{Research questions left by the implementation}

The most valuable unresolved questions are practical rather than foundational.
How often is the target-generated ideal small in real Lean developments? When
does fixed-degree affine lifting beat direct ideal closure? Which generator
orders minimize kernel replay rather than search alone? Can the binomial
certificate basis be enlarged while keeping endpoints polynomial? Which exact
Leant candidate grammars admit useful observational congruences?

Each has a concrete experiment. Compare discovered generator counts under
alternative term orders while replaying the same source goals. Measure
monomial-lift dimensions and polynomial certificate sizes on matched affine
problems. Compare proof construction and kernel time for expanded versus shared
multiplier expressions. Add higher-shift binomial factors and require support
and pole tests before admitting the new template. Restrict semantic deduplication
to a source grammar with proved congruence lemmas and retain a no-pruning control.

None requires treating an unproved CAS assertion as evidence. The architecture
can remain the one Forge already has; the advance comes from choosing
mathematical fragments with finite, inspectable certificates.
